% Value-Aware Token Pruning (VATP) for KV Cache Reduction
% Authors: Zhiyu Guo, Hidetaka Kamigaito, Taro Watanabe (NAIST)
% Consolidated technical content

\section{Abstract}
Scaling context size of LLMs enables new tasks but KV cache memory significantly limits practical applications. Recent token pruning methods for KV cache reduction rely solely on attention scores. However, value vector norms are notably non-uniform. We propose Value-Aware Token Pruning (VATP) using both attention scores and L1 norm of value vectors to evaluate token importance. VATP outperforms attention-score-only baselines (H2O, Scissorhands) in 12-14 out of 16 LongBench tasks at 50% KV cache budget.

\section{Key Observations}

\subsection{Attention Output Formula}
\begin{equation}
\text{Attention}(Q, K, V)_t = \sum_{i \leq t} a_i^t \mathbf{v}_i
\end{equation}
where $a_i^t$ is attention score of query token $t$ to token $i$, and $\mathbf{v}_i$ is value state.

Each token's influence on attention output is determined by BOTH:
1. Attention score $a_i^t$
2. Value vector $\mathbf{v}_i$

\subsection{Critical Finding: Value Vector Norm Distribution}
The L1 norm of value vectors is highly non-uniformly distributed across all layers and heads.

\textbf{Attention Sink Phenomenon:}
- First two tokens (start token and first period/newline) have massive attention scores
- BUT their value vector norms are close to 0
- This means despite huge attention scores, their actual contribution to output is minimal

\subsection{Opposite Case}
In some heads of the last layer:
- Second attention sink token has SMALLER attention score than other tokens
- BUT its L1 norm is significantly LARGER than others

\section{Methodology}

\subsection{Baseline Attention Score Methods}
\textbf{H2O (Accumulated):}
\begin{equation}
S_k^t = \sum_{k \leq j \leq t} a_k^j
\end{equation}

\textbf{Scissorhands (History Window):}
\begin{equation}
S_k^t = \sum_{\max(t-w, k) \leq j \leq t} a_k^j
\end{equation}

\subsection{VATP Token Importance Metric}
Product of attention score and value vector norm:
\begin{equation}
I_k^t = S_k^t \cdot \|\mathbf{v}_k\|_1
\end{equation}

where $\|\mathbf{v}_k\|_1$ is L1 norm of token $k$'s value vector.

\subsection{Handling Attention Sinks}
Since attention sink tokens have very small L1 norms, VATP would accidentally remove them. But removing them shifts attention distribution dramatically. Solution: intentionally keep first F tokens (F=20 for LLaMA2-7B, F=40 for Vicuna).

\section{Experiments}

\subsection{Setup}
- Models: LLaMA2-7B-chat, Vicuna-v1.5-7B
- Benchmark: 16 English tasks from LongBench
- KV cache budget: 50% of input prompt length
- Tasks: Single/Multi-Doc QA, Summarization, Few-shot, Synthetic, Code

\subsection{Main Results}
\textbf{LLaMA2-7B-chat:}
- VATP surpasses H2O in 12/16 tasks
- VATP surpasses Scissorhands in 13/16 tasks

\textbf{Vicuna-v1.5-7B:}
- VATP exceeds H2O in 12/16 tasks
- VATP exceeds Scissorhands in 14/16 tasks

Key: Where VATP doesn't surpass baseline, it's very comparable. But in some tasks (e.g., MultiFieldQA), VATP significantly outperforms.

\subsection{KV Budget Ratio Variation}
At high reduction ratios, Scissorhands w/ VATP significantly outperforms plain Scissorhands.

\subsection{Inference Efficiency}
- VATP adds negligible overhead
- Value vector norm size = 1/(2*d_head) of KV cache size
- No significant difference in throughput or memory vs baseline
- Compatible with FlashAttention when using Scissorhands

\subsection{Throughput Improvement}
At 50% KV budget: 1.53-1.58x speedup
At 25% KV budget: 2.08-2.26x speedup

\subsection{Ablation: Norm Type}
L1 norm performs best (average F1: 28.00) vs L2 (27.70) and L-infinity (27.71)

\section{Conclusions}
1. Value vector norms are critically important for token importance evaluation
2. Attention score alone is insufficient
3. VATP: simple product of attention score and L1 value norm
4. Negligible computational overhead, no fine-tuning required
5. Challenges prevailing belief that "attention score is all you need"

\section{Limitations}
1. H2O integration requires attention matrix materialization (incompatible with FlashAttention)
2. Not directly applicable to Grouped-Query Attention (GQA)
3. Future: pruning based on both key AND value vector norms
